%%%%%%%% AI4Physics @ ICML 2026 — clax paper %%%%%%%%%%%%%%%%%

\documentclass{article}

\usepackage{microtype}
\usepackage{graphicx}
\usepackage{subcaption}
\usepackage{booktabs}
\usepackage{hyperref}
\newcommand{\theHalgorithm}{\arabic{algorithm}}
\usepackage{icml2026}
\usepackage{amsmath}
\usepackage{amssymb}
\usepackage{mathtools}
\usepackage[capitalize,noabbrev]{cleveref}
\usepackage{xspace}

% Custom commands
\newcommand{\clax}{\textsc{clax}\xspace}
\newcommand{\claxpt}{\textsc{clax-pt}\xspace}
\newcommand{\class}{\textsc{class}\xspace}
\newcommand{\classpt}{\textsc{class-pt}\xspace}
\newcommand{\discoeb}{\textsc{disco-eb}\xspace}
\newcommand{\abcmb}{\textsc{abcmb}\xspace}
\newcommand{\jax}{\textsc{jax}\xspace}
\newcommand{\Cl}{C_\ell}
\newcommand{\Pk}{P(k)}

\icmltitlerunning{From Einstein-Boltzmann Equations to Galaxy Power Spectra}

\begin{document}

\twocolumn[
  \icmltitle{From Einstein-Boltzmann Equations to Galaxy Power Spectra: \\
    Physicist-Supervised, Oracle-Guided AI Reimplementation in JAX}

  \icmlsetsymbol{equal}{*}

  \begin{icmlauthorlist}
    \icmlauthor{Siddharth Mishra-Sharma}{anthropic}
    \icmlauthor{Minh Nguyen}{kavli}
  \end{icmlauthorlist}

  \icmlaffiliation{anthropic}{Anthropic, San Francisco, CA, USA}
  \icmlaffiliation{kavli}{Kavli IPMU, University of Tokyo, Kashiwa, Japan}

  \icmlcorrespondingauthor{Minh Nguyen}{nhat.minh.nguyen.111@gmail.com}

  \icmlkeywords{differentiable programming, cosmology, Boltzmann solver, perturbation theory, AI-assisted development}

  \vskip 0.3in
]

\printAffiliationsAndNotice{}

% ===========================================================================
\begin{abstract}
Bayesian inverse problems in cosmology are tackling increasingly high dimensional parameter spaces, driven by the need to marginalize over complex astrophysical uncertainties and instrumental systematics.
Very high-dimensional inference problems require many (expensive) simulations and remain challenging for simulation-based inference methods such as normalizing flows and diffusion models. Sampling-based methods---typically relying on gradient-based Monte Carlo Markov chain (MCMC) method, e.g., Hamiltonian Monte Carlo (HMC)---therefore remain the standard approach. HMC requires access to the gradient of the posterior with respect to model input parameters, hence differentiable theory codes. Meanwhile, cosmological analyses are also moving towards joint analyses of multiple data sets, e.g., those of the cosmic microwave background (CMB) and galaxy redshift surveys. However, no existing code provides both the full set of CMB power spectra and the galaxy power spectra with exact automatic differentiation on GPU.
We address this gap with \clax and \claxpt: fully differentiable \jax reimplementations of the \class Boltzmann solver and its one-loop perturbation theory extension \classpt---developed via physicist-supervised, oracle-guided AI agents that iterate against reference data and source codes from \class and \classpt---autonomously resolving implementation bugs and escalating to human review when oracle feedback stalls.
Across the development, agents autonomously resolved 38 of 40 documented bugs (95\%) without human input.
The remaining 5\% required human physics judgment and fell into exactly two failure modes: \emph{architectural incoherence} (code structure incompatible with the target physics) and \emph{oracle-passing fudge factors} (a numerical patch satisfying all tests but encoding no physical model).
\clax achieves $< 0.2\%$ lensed CMB accuracy at $\ell \leq 2000$ and $< 0.6\%$ unlensed $\Cl^{TT}$ at $\ell \leq 1000$; \claxpt achieves $< 0.7\%$ galaxy power spectrum accuracy for monopole and quadrupole spectra and $< 1.5\%$ for hexadecapoles at $k < 0.3\;h/$Mpc.
Automatic differentiation gradients agree with fourth-order finite differences to $0.15\%$ on average.
The full pipeline runs at 34\,s/step on a single V100 GPU with HMC-ready accuracy.
We characterize the boundary between unsupervised and supervised modes of AI-assisted scientific development, propose three principles for structuring human--agent collaboration, and compare against concurrent differentiable cosmology codes \discoeb and \abcmb.
\end{abstract}

% ===========================================================================
\section{Introduction}
\label{sec:intro}

The standard model of cosmology is tested by comparing theoretical predictions against observations from galaxy surveys such as DESI \citep{desi2025} and Euclid \citep{euclid2024}, and CMB experiments including Planck \citep{planck2020}.
This comparison increasingly benefits from gradient-based inference: exact Fisher matrices require derivatives $\partial \Cl / \partial \theta_i$, HMC sampling \citep{neal2011} requires $\nabla_\theta \log P(\boldsymbol{\theta} | \mathrm{data})$, and even simulation-based inference \citep{cranmer2020}---which does not require gradients---can leverage them for score-matching and summary-network training.
Gradient-based inference libraries such as NumPyro \citep{phan2019}, BlackJAX \citep{cabezas2024}, and normalizing-flow-augmented MCMC \citep{gabrie2022} are mature, but the missing ingredient is a differentiable cosmological forward model that can be passed to them.

The dominant Boltzmann codes---\class \citep{blas2011} and CAMB \citep{lewis2000}---are written in C and Fortran, are not differentiable, and are not GPU-native.
Several \jax implementations exist but cover only subsets of the problem: \discoeb \citep{hahn2023} solves the linear Einstein--Boltzmann equations to compute the matter power spectrum $\Pk$ but not CMB spectra or galaxy clustering observables; \abcmb \citep{zhou2026} implements the CMB pipeline including lensing and massive neutrinos but does not include perturbation theory for galaxy surveys; \textsc{jax-cosmo} \citep{lemos2023} computes the nonlinear matter power spectrum via fitting formulae but does not solve the full Boltzmann hierarchy; and SymBoltz.jl \citep{sletmoen2026} provides a full Boltzmann solver in Julia but without GPU support.
No existing code provides the full chain from cosmological parameters through CMB angular power spectra \emph{and} one-loop EFT galaxy power spectrum multipoles, with end-to-end automatic differentiation on GPU.

We describe the development of \clax and \claxpt---fully differentiable \jax reimplementations of the \class \citep{blas2011} and \classpt \citep{chudaykin2021} pipelines---and the oracle-guided agentic methodology used to build them.
Using \class v3.3.4 and \classpt as ground truth, we find a clear boundary: agents handle implementation bugs (convention errors, algorithmic transcription, coefficient verification) autonomously and reliably, but fail at two classes of problem that require human physics judgment: (1) recognizing that an architecture is structurally incompatible with the intended physics, and (2) detecting that a numerically effective patch satisfies all oracle tests at the calibration point while encoding incorrect physics.

\paragraph{Contributions.}
\begin{itemize}
    \item \clax: a \jax reimplementation of the full \class CMB pipeline (background, thermodynamics, perturbations, transfer, lensing), achieving $< 0.1\%$ accuracy at $\ell \leq 2000$ with JIT-compiled GPU execution at 34\,s/step (V100, \texttt{fit\_cl} preset).
    \item \claxpt: a \jax reimplementation of \classpt's one-loop EFT power spectra with IR resummation and RSD multipoles, achieving $< 0.7\%$ accuracy for monopole and quadrupole spectra.
    \item Validated autodifferentiation gradients matching finite differences to $0.15\%$ across cosmological parameters.
    \item A landscape comparison against \discoeb, \abcmb, \class, and \classpt, showing that \clax is the first code combining CMB and one-loop galaxy spectra with full AD on GPU.
    \item An empirical bug taxonomy of 40 bugs across the development, classifying autonomous vs.\ human-required resolution and identifying two specific failure modes of oracle-guided development.
\end{itemize}

% ===========================================================================
\section{Background}
\label{sec:background}

\subsection{Cosmological inference and differentiable Boltzmann codes}

Galaxy surveys and CMB experiments measure the angular power spectrum $\Cl$ and the galaxy power spectrum $P_\mathrm{gg}(k, \ell)$, which depend on cosmological parameters $\boldsymbol{\theta}$ (density, dark matter density, Hubble constant, spectral index, etc.)\ through the Boltzmann equation---a coupled system of ODEs for photon, baryon, dark matter, and neutrino perturbations that must be integrated numerically.
Both are expensive to approximate by finite differences and practical to compute exactly only with automatic differentiation (AD) for codes of this complexity.
While concurrent work (\abcmb, SymBoltz.jl) provides AD for CMB spectra, \clax is the first code to combine differentiable CMB $\Cl$ with one-loop galaxy power spectrum multipoles, enabling \texttt{jax.grad} through the full cosmological analysis pipeline for galaxy clustering.

\subsection{AI-assisted scientific code development}

The oracle-based methodology we employ is inspired by the AI C compiler project \citep{carlini2026}, in which parallel Claude Code agents produced a Rust C compiler using GCC as an oracle.
Key infrastructure includes: (i) a shared \texttt{CHANGELOG.md} updated after every session to prevent re-exploration of dead ends; (ii) test-first development where every module's tests are written against reference data before implementation; (iii) a \texttt{--fast} flag sampling 10\% of test cases to prevent ``time blindness''; and (iv) parallel agent teams using git worktrees to explore competing hypotheses independently.

The critical challenge not addressed in the compiler setting is \emph{domain knowledge}: \class is not just a large codebase---it encodes decades of cosmological physics knowledge, conventions, and approximations that require physics expertise to interpret.
This makes the boundary between supervised and unsupervised development especially important to characterize.

% ===========================================================================
\section{Methodology}
\label{sec:methods}

\subsection{Pipeline architecture}

\clax implements a sequential pipeline of pure functions, each returning a frozen PyTree:
\begin{equation*}
\text{CosmoParams} \to \text{BG} \to \text{TH} \to \text{EB} \to \text{PR} \to \text{TR} \to \text{HR} \to \text{LE}
\end{equation*}
where BG is background cosmology, TH is thermodynamics/recombination, EB is the Einstein--Boltzmann perturbation hierarchy, PR is the primordial spectrum, TR is transfer functions, HR is harmonic ($\Cl$) integration, and LE is CMB lensing via Wigner $d$-functions.
\texttt{CosmoParams} fields are \jax-traced (for AD); \texttt{PrecisionParams} fields are static (control array shapes).
No branching on traced values ensures compatibility with \texttt{jax.jit} and \texttt{jax.grad}.

The Einstein--Boltzmann system is integrated using Kvaerno5 (ESDIRK, default) or Rodas5 (Rosenbrock) solvers from the Diffrax library \citep{kidger2021}, with a DISCO-EB-style filtered PID step controller using $k$-dependent weighting on six key perturbation variables.
Reverse-mode AD through the ODE solve uses \texttt{RecursiveCheckpointAdjoint} for memory efficiency, with \texttt{custom\_vjp} wrappers for the shooting method ($\theta_s \to H_0$ via implicit differentiation) and the Saha equilibrium (implicit-function-theorem JVP).

\claxpt extends the pipeline with one-loop EFT power spectra following the \classpt algorithm: FFTLog decomposition of $P_{22}$ and $P_{13}$ integrals \citep{fang2017}, precomputed $M_{22}/M_{13}$ kernel matrices, IR resummation via DST with odd/even spline mode removal \citep{senatore2015, vlah2015}, an EFT bias expansion ($b_1, b_2, b_{s^2}, b_{3\mathrm{nl}}$, counterterms, stochastic contributions), and RSD multipoles ($\ell = 0, 2, 4$) via Gauss--Legendre quadrature over $\mu$ with anisotropic BAO damping.
The entire \claxpt module ($\sim$2,100 lines) is end-to-end differentiable: \texttt{jax.grad} flows through FFTLog, IR resummation, bias assembly, and GL quadrature.

\subsection{Development methodology}

The \clax CMB pipeline ($\sim$10,000 lines) was developed over $\sim$6 days across $\geq$100 Claude Code (Opus 4.6) agent sessions.
The \claxpt one-loop PT module ($\sim$2,100 lines) was developed over 12 days and 57 logged sessions.
Both phases used the same oracle-guided methodology: every module has a test file written \emph{before} implementation, specifying what \class (or \classpt) produces; the code is finished when and only when the oracle tests pass.

The supervision methodology---encoded in \texttt{CLAUDE.md} and refined across both development phases---proved essential.
Two rules were particularly important: (1) never add fudge factors (if a test fails at 0.2\%, find the actual bug), and (2) test at multiple parameter points, not just the fiducial cosmology.
Both rules directly prevented specific failure modes documented in \cref{sec:supervised}.

% ===========================================================================
\section{Results: Accuracy and Performance}
\label{sec:results}

\subsection{Landscape comparison}

\Cref{tab:landscape} compares \clax against existing differentiable and non-differentiable cosmology codes.
\clax is the first \jax code providing both CMB angular power spectra and one-loop galaxy power spectrum multipoles with end-to-end automatic differentiation on GPU.

\begin{table*}[t]
\caption{Comparison of cosmological theory codes. ``AD'' = automatic differentiation. ``PT'' = one-loop perturbation theory for galaxy clustering. \clax is the only code combining CMB $\Cl$, lensing, $\Pk$, and 1-loop PT with full AD on GPU.}
\label{tab:landscape}
\centering
\small
\begin{tabular}{lccccccc}
\toprule
Code & Framework & $\Cl$ & Lensing & $\Pk$ & 1-loop PT & AD & GPU \\
\midrule
\class + \classpt & C / Python & \checkmark & \checkmark & \checkmark & \checkmark & --- & --- \\
CAMB & Fortran & \checkmark & \checkmark & \checkmark & --- & --- & --- \\
\discoeb & \jax & --- & --- & \checkmark & --- & \checkmark & \checkmark \\
\abcmb & \jax & \checkmark & \checkmark & \checkmark & --- & \checkmark & \checkmark \\
SymBoltz.jl & Julia & \checkmark & --- & \checkmark & --- & \checkmark & --- \\
\textbf{\clax + \claxpt} & \jax & \checkmark & \checkmark & \checkmark & \checkmark & \checkmark & \checkmark \\
\bottomrule
\end{tabular}
\end{table*}

\subsection{CMB pipeline accuracy}

\Cref{tab:cmb_accuracy} summarizes accuracy of \clax against \class at the Planck 2018 best-fit cosmology (\texttt{planck\_cl} preset).
Background quantities match to $< 0.01\%$; thermodynamics to $< 0.1\%$.
CMB temperature and polarization spectra are within $0.02\%$--$0.57\%$ across $\ell = 20$--$1000$.
Lensed spectra at $\ell = 2000$ pass at $< 0.2\%$.
Multi-cosmology validation across 10 parameter points (omega\_b/cdm $\pm 20\%$, $h \pm 10\%$, massive neutrinos 0.06--0.3~eV, $N_\mathrm{eff} = 2.0$--$4.0$) confirms sub-0.5\% accuracy throughout parameter space.

\begin{table}[t]
\caption{CMB accuracy (\clax vs.\ \class, \texttt{planck\_cl} preset, Planck 2018 fiducial).}
\label{tab:cmb_accuracy}
\centering
\small
\begin{tabular}{ll}
\toprule
Quantity & Max error \\
\midrule
$H(z)$, $D_A(z)$, $r_s$ & $< 0.01\%$ \\
$x_e(z)$, $g(\tau)$ & $< 0.1\%$ \\
$\Cl^{TT}$, $\ell = 20, 100, 500, 1000$ & 0.08\%, 0.02\%, 0.15\%, 0.57\% \\
$\Cl^{EE}$, $\ell = 20, 100, 500, 1000$ & 0.21\%, 0.02\%, 0.15\%, 0.26\% \\
Lensed $\Cl$, $\ell = 2000$ & $< 0.2\%$ \\
Multi-cosmology (10 pts) & $< 0.5\%$ \\
\bottomrule
\end{tabular}
\end{table}

\subsection{Galaxy power spectrum accuracy}

\Cref{tab:pt_accuracy} shows accuracy for one-loop EFT galaxy power spectra at $z = 0.38$ (BOSS effective redshift) against \classpt.
Real-space spectra ($P_{mm}$, $P_{gg}$, $P_{gm}$) agree to $< 0.31\%$ (mean 0.04\%).
RSD monopoles ($\ell = 0$) agree to $< 0.6\%$ for matter and galaxies.
Quadrupoles ($\ell = 2$) agree to $< 0.7\%$ for matter and $< 0.9\%$ for galaxies.
Hexadecapoles ($\ell = 4$) are validated using $|\Delta|/\max(|\mathrm{ref}|) < 2\%$ since \classpt values are small and the spectrum crosses near zero.

\begin{table}[t]
\caption{Galaxy power spectrum accuracy (\claxpt vs.\ \classpt, $z = 0.38$).}
\label{tab:pt_accuracy}
\centering
\small
\begin{tabular}{lll}
\toprule
Spectrum & Max error & Metric \\
\midrule
$P_{mm}^\mathrm{real}$, $P_{gg}^\mathrm{real}$, $P_{gm}^\mathrm{real}$ & 0.31\% & rel. \\
$P_{mm}^{(\ell=0)}$, $P_{gg}^{(\ell=0)}$ & 0.59\%, 0.56\% & rel. \\
$P_{mm}^{(\ell=2)}$, $P_{gg}^{(\ell=2)}$ & 0.70\%, 0.89\% & rel. \\
$P_{mm}^{(\ell=4)}$, $P_{gg}^{(\ell=4)}$ & 0.70\%, 1.43\% & abs/max \\
\bottomrule
\end{tabular}
\end{table}

\subsection{Gradient verification}

Automatic differentiation gradients (reverse-mode, using \texttt{jax.grad}) agree with fourth-order centered finite differences to $0.15\%$ on average, with 97.7\% of $k$-modes passing a 1\% threshold.
\Cref{tab:gradients} shows the parameter-by-parameter breakdown for the matter power spectrum $P_{mm}(k)$ at $k = 0.1\;h/$Mpc.

\begin{table}[t]
\caption{AD vs.\ finite-difference gradients for $P_{mm}(k=0.1\;h/\mathrm{Mpc})$.}
\label{tab:gradients}
\centering
\small
\begin{tabular}{lcc}
\toprule
Parameter & $|\mathrm{AD} / \mathrm{FD} - 1|$ & Status \\
\midrule
$\ln(10^{10} A_s)$ & $< 0.01\%$ & \checkmark \\
$n_s$ & $< 0.01\%$ & \checkmark \\
$h$ & $< 0.5\%$ & \checkmark \\
$\omega_b$ & $< 1.0\%$ & \checkmark \\
$\omega_\mathrm{cdm}$ & $< 0.5\%$ & \checkmark \\
\bottomrule
\end{tabular}
\end{table}

Key technical ingredients enabling stable AD through the full pipeline include: (1) \texttt{RecursiveCheckpointAdjoint} for memory-efficient reverse-mode through the stiff perturbation ODE; (2) a \texttt{custom\_vjp} on the shooting method ($\theta_s \to H_0$) implementing implicit differentiation; and (3) a \texttt{custom\_jvp} on the hydrogen Saha equilibrium equation, replacing explicit \texttt{stop\_gradient} calls with implicit-function-theorem tangents to restore nonzero sensitivities around recombination.

\subsection{Performance}

\clax offers two presets for different use cases (\cref{tab:performance}).
The \texttt{fit\_cl} preset uses 20~$k$/decade, $\ell_\mathrm{max} = 17$ hierarchy, table-based Bessel functions, and $\texttt{rtol} = 10^{-3}$ to achieve 34\,s/step on a V100 GPU---sufficient for HMC.
The \texttt{planck\_cl} preset uses 60~$k$/decade, $\ell_\mathrm{max} = 50$, and tight tolerances for science-grade accuracy at $\sim$487\,s on H100.
An alternative Rodas5 Rosenbrock solver \citep{dimarzio1993}, following \discoeb, avoids Newton iteration entirely (one LU factorization per step instead of iterative convergence).
A batched variant (\texttt{Rodas5Batched}) solves all $k$-modes in a single \texttt{diffeqsolve} call with shared adaptive time-stepping across the batch, vmapping the Jacobian computation and LU factorization---the same architecture used by \discoeb.
On CPU, the unbatched Rodas5 yields a measured $\sim$1.3$\times$ speedup over Kvaerno5; the batched variant is expected to offer substantially larger gains on GPU by saturating GPU throughput with large batched linear algebra.

\begin{table}[t]
\caption{Performance summary.}
\label{tab:performance}
\centering
\small
\begin{tabular}{lccl}
\toprule
Preset & GPU & Time & $\Cl$ accuracy \\
\midrule
\texttt{fit\_cl} & V100-32GB & 34\,s & $<1.5\%$ ($\ell \leq 500$) \\
\texttt{planck\_cl} & H100-80GB & 487\,s & $<0.2\%$ ($\ell \leq 2000$) \\
\bottomrule
\end{tabular}
\end{table}

\subsection{Application: gradient scaling}

The primary advantage of AD over finite differences for cosmological inference is the scaling of gradient cost with the number of parameters $d$.
Finite-difference gradients require $2d$ forward evaluations (centered differences); reverse-mode AD requires one forward pass plus one backward pass regardless of $d$.
The backward pass typically costs 2--4$\times$ the forward pass, so the effective speedup for $d = 6$ standard $\Lambda$CDM parameters is $2d / (1 + c_\mathrm{bwd}) \approx 3$--$4\times$, growing to $\sim$10$\times$ for extended models with $d = 10$--$20$ parameters (neutrino masses, dark energy, bias parameters, counterterms).
The key advantage is the asymptotic scaling: AD gradient cost is $O(1)$ in $d$, while finite differences scale as $O(d)$.
Combined with the 34\,s forward evaluation of \texttt{fit\_cl}, this makes gradient-based sampling with exact first-principles theory predictions practical for full-shape galaxy clustering analysis.

% ===========================================================================
\section{Supervised vs.\ Unsupervised Development}
\label{sec:supervised}

\subsection{Bug taxonomy}

Over the full development of \clax and \claxpt, we documented 40 bugs: 26 in the CMB pipeline and 14 in the PT module.
We classify these into four types by their nature and whether autonomous resolution was possible.

\textbf{Type~I: Convention and unit bugs.}
Errors in translating \class conventions into \jax: wrong sign, missing numerical factor, incorrect unit conversion.
\emph{Agents resolved all Type~I bugs autonomously} by reading the corresponding \class source code and correcting the factor.

\textbf{Type~II: Algorithm implementation bugs.}
Errors in the numerical algorithm itself: wrong grid type, incorrect matrix storage format, incomplete angular kernel functions.
\emph{Agents resolved all Type~II bugs autonomously}, typically by careful comparison of intermediate arrays against \class/\classpt at specific $(k, z)$ points.

\textbf{Type~III: Architectural bugs.}
The computational structure assumed by the implementation is fundamentally incapable of representing the correct physics.
One instance occurred: the analytic Legendre projection architecture assumed isotropic BAO damping, making it unable to handle anisotropic damping $\Sigma_\mathrm{tot}^2(\mu)$.
\emph{This required human intervention.}

\textbf{Type~IV: Physics judgment failures.}
A numerical patch satisfies all oracle tests but encodes physically unmotivated behavior.
One instance occurred: a scalar fudge factor $\alpha = 0.27$ passed all fiducial tests but has no physical derivation.
\emph{This required human intervention.}

The resolution breakdown is: Types~I+II (38 bugs): 100\% autonomous; Types~III+IV (2 bugs): 0\% autonomous.
Overall, 38/40 bugs (95\%) were resolved without human input.

\subsection{Principles for human--agent collaboration}

We distill three principles from the empirical data:

\textbf{P1: Oracle testing verifies \emph{what}, not \emph{why}.}
An oracle compares numerical values at a reference parameter point.
It catches implementation errors (wrong sign, wrong coefficient) but not physics errors (wrong model, numerically compensated).
A fudge factor with zero degrees of freedom at the calibration point will pass every oracle test at that point.
\emph{Defense}: test across diverse parameter space; test intermediate quantities; build auxiliary tests that check structural properties (e.g., that $P_\mathrm{tree}(k,\mu)$ has the correct anisotropy as a function of $\mu$), not just numerical values.

\textbf{P2: Shared memory prevents re-exploration but not architectural loops.}
\texttt{CHANGELOG.md} records which algorithms were tried and failed, preventing agents from rediscovering dead ends.
This was effective for Type~I and~II bugs.
It did not prevent 33 sessions of exploration within the wrong architectural framework because all attempts were coherent within the (incorrect) model.
\emph{Defense}: when a bug persists after $N \approx 5$--10 sessions without progress, treat the lack of progress itself as evidence of an architectural problem and trigger human review.

\textbf{P3: The irreducible human role is architectural and physical judgment.}
Agents excel at reading \class source code and transcribing equations faithfully; identifying which term disagrees via targeted comparison with oracle outputs; and adjusting coefficients and conventions.
Agents fail at recognizing that an architectural assumption is inconsistent with the physics, and at detecting that a numerically adequate patch produces zero test failures but is not physics.
\emph{Detecting physics errors that pass all tests requires knowledge of what the physics should look like}, which is a form of human supervision that oracle testing cannot replace.

% ===========================================================================
\section{Related Work}
\label{sec:related}

\textbf{Differentiable cosmological codes.}
\discoeb \citep{hahn2023} solves the linear Einstein--Boltzmann equations in \jax to compute the matter power spectrum but does not compute CMB spectra ($\Cl$) or galaxy clustering observables.
\abcmb \citep{zhou2026} provides a complete CMB pipeline in \jax including lensing, massive neutrinos, and a state-of-the-art HyRex recombination treatment, but does not include perturbation theory for galaxy surveys.
SymBoltz.jl \citep{sletmoen2026} is the closest comparison in design philosophy: it is approximation-free (no TCA/RSA/UFA switching), uses Rodas5P (the same Rosenbrock solver family as \clax), achieves 0.1\% accuracy against \class, and supports forward-mode AD for Fisher forecasts.
However, SymBoltz.jl operates in Julia without GPU support, and its reverse-mode AD is not yet fast enough for MCMC with perturbation-derived spectra \citep{sletmoen2026}.
\textsc{jax-cosmo} \citep{lemos2023} computes the nonlinear matter power spectrum via fitting formulae but does not solve the full Boltzmann hierarchy.
\clax is the first \jax code to provide full CMB spectra plus one-loop galaxy power spectra with exact reverse-mode automatic differentiation and GPU execution.

\textbf{EFT power spectra and IR resummation.}
The EFT of large-scale structure \citep{perko2016, damico2020} provides a systematic expansion for the galaxy power spectrum.
\classpt \citep{chudaykin2021} is the standard implementation; \claxpt reproduces it in differentiable \jax.
IR resummation handles BAO peak broadening from long-wavelength displacements \citep{senatore2015, vlah2015}.

\textbf{Agent-assisted scientific software.}
Recent work has shown that LLM agents can complete extended scientific tasks autonomously \citep{carlini2026}, but most evaluations focus on syntactic correctness (compilation, test passage) rather than physical validity.
Our work identifies semantic correctness---the requirement that a passing numerical patch has physical basis, not just calibration-point accuracy---as the key challenge distinguishing scientific software from general code generation.

% ===========================================================================
\section{Conclusion}
\label{sec:conclusion}

We have developed \clax and \claxpt---fully differentiable \jax reimplementations of the \class and \classpt Boltzmann codes---using oracle-guided agentic development.
The codes achieve $< 0.2\%$ lensed CMB accuracy at $\ell \leq 2000$, $< 0.7\%$ galaxy power spectrum accuracy for monopole and quadrupole spectra, with automatic differentiation gradients verified to $0.15\%$ and GPU execution at 34\,s/step on V100.
Empirically, AI agents resolved 95\% of implementation bugs autonomously.
The remaining 5\% required human physics judgment: recognizing a structural mismatch between an architecture and its target physics, and detecting a numerically effective but physically unmotivated fudge factor.

These findings suggest a practical framework for AI-assisted scientific code development: agents operate autonomously in unsupervised mode for implementation work, but human scientists remain in the loop as architectural and physical reviewers, intervening proactively when oracle feedback stalls or when a numerical patch lacks physical derivation.
The codes are open source at \url{https://github.com/smsharma/clax}.

% ===========================================================================
\bibliography{references}
\bibliographystyle{icml2026}

\end{document}
